\begin{proof}
Write $t=s^2$ and let $N=|\Delta|n+1$.  Every feasible $V$ has the
exact compact parametrization
\begin{equation}
 V=[sz,W_s(z)Z],\qquad
 \Rnorm z=1,\quad Z^*Z\preceq I_{n-1},\qquad
 W_s(z)=I_N-(1-\sqrt{1-s^2})zz^*.
 \label{I:R:eq:terminal-parametrization}
\end{equation}
Indeed this is the Schur complement of $I-V^*V$.  For a polynomial
$\tau$, the quotient $T_\tau e_1/s$ has an analytic extension to
$s=0$.  For general $\tau\in\ell^2$, the tail bound from
\eqref{I:R:eq:terminal-refined-energy}, divided by $|s|$, is at most
the terminal tail's $\ell^2$ norm.  Polynomial approximation therefore
gives a uniform continuous extension on the compact parameter domain.

At $s=0$ write
\[
 B_j(0)=\begin{pmatrix}0&b_j\\0&C_j\end{pmatrix},\qquad
 \ell(0)=(0,l),\qquad z_j=(\alpha_j,\xi_j),
\]
where $z_j$ is the first-column derivative block for $B_j$, and let
$z_o$ be its scalar output block.  The limiting divided terminal row is
\begin{equation}
 \overline{\tau(\varepsilon)}z_o+
 \sum_j R_j(C,l)\xi_j,
 \qquad
 R_j(C,l)=\sum_v\overline{\tau(jv)}\,lC_v.
 \label{I:R:eq:terminal-divided-limit}
\end{equation}
The $\alpha_j$ do not contribute.  Maximizing over unit $z$ gives
$|\tau(\varepsilon)|^2+\sum_j\Rnorm{R_j(C,l)}^2$.
The remaining columns in \eqref{I:R:eq:terminal-parametrization} imply
$\sum_j C_j^*C_j+l^*l\preceq I_{n-1}$.  Conversely every such reduced
tuple is feasible by choosing $b_j=0$.  Uniform convergence of the
compact objective proves
\[
 a_{\tau,n}=|\tau(\varepsilon)|^2+
 \max_{\sum C_j^*C_j+l^*l\preceq I_{n-1}}
                    \sum_j\Rnorm{R_j(C,l)}^2.
\]
For $n=1$ the reduced state space is zero and the sum is zero.

This passive maximum equals the invariant-space maximum in
\eqref{I:R:eq:terminal-leading-coefficient}.  A canonical observation map
on $K$ is its isometric inclusion, so $R_j$ has norm
$\Rnorm{P_K\tau_j}$.  Conversely a passive reduced tuple has a
contractive observation map $O$ whose finite-dimensional range $K$ is
backward invariant.  Since $OO^*\preceq P_K$,
$\Rnorm{O^*\tau_j}\le\Rnorm{P_K\tau_j}$.  This gives the reverse
inequality after summing.  A maximizing passive tuple exists by
compactness and uniform convergence; its range attains the
invariant-space maximum as well.

Finally write \eqref{I:R:eq:terminal-elimination} as
$P_{\tau,n}(q)=\max_t K_q(t)h_{\tau,n}(t)/t$.
The trial $t=q$ tends to $a_{\tau,n}$.  On $t\ge\eta>0$ the kernel
$K_q(t)$ tends uniformly to zero, whereas on $0<t<\eta$ it is at most
one and $h(t)/t\to a_{\tau,n}$.  Splitting at $\eta$ and then sending
$\eta$ to zero proves the upper limit.  This also covers
$a_{\tau,n}=0$.
\end{proof}
